\section{Introduction}
Quantitative research reports distribute information across articles, tables, figures, supplements, protocols, and statistical analysis plans. Readers must reconcile these sources to identify the population, outcome, analysis, and uncertainty represented by a result. This task also precedes evidence synthesis: reliable extraction depends on interpretable source reports. Xu et al. found that extraction errors in systematic reviews of adverse events sometimes changed the direction or statistical significance of pooled effects.\cite{xu2022extraction} Checking source reporting therefore complements efforts to improve extraction and synthesis.

Our project arose from peer-review experience with inconsistent statistical results across manuscript components. An earlier feasibility exercise applied a large language model (LLM) to one JAMA randomized trial and its supplements, producing candidate concerns in the supplements for further inspection. This experience motivated a broader question: can a structured workflow identify specific reporting concerns in published clinical studies while preserving the evidence needed for human assessment?

LLMs have been applied to statistical analysis, reference generation, and clinical triage. Evaluations by Ruta et al., Graf et al., and Paris et al. illustrate their potential uses alongside dependence on instructions, inaccurate references, and limited agreement with specialists.\cite{ruta2025statistics,graf2024referencing,paris2026triage} Document-checking studies address a more specific problem. The FIND, FLAWS, and SPOT preprints evaluate recovery of inserted inconsistencies, inserted scientific errors, and documented errors associated with corrections or retractions, respectively.\cite{lovering2025find,xi2025flaws,son2025spot} Bianchi et al. instead obtained human assessments of newly proposed findings in artificial intelligence papers.\cite{bianchi2025errors}

In ophthalmology, La, Paris, and colleagues assessed detection, localization, explanation, and correction of planted inconsistencies and reviewed additional findings.\cite{la2026consistency} Dong et al. evaluated structured checking of categorical analyses, identifying useful capabilities alongside extraction failures and omitted analyses.\cite{dong2026statistical} These studies motivate supplementary review and distinguish recovery of known errors from assessment of newly generated concerns. Their different targets and reference standards limit direct performance comparisons.

A practical workflow must also record what was examined and how each concern arose. Persistent source maps, numerical relationships, and screening decisions can preserve context across files and review stages. Document tools, reusable extraction assets, and delegated tasks provide a way to organize this work, provided that derived evidence remains linked to original sources. Human responsibility and transparent reporting of AI use remain necessary.\cite{flanagin2026guidance}

We present TRACE-Q, an exploratory multi-agent toolkit for evidence-grounded review of quantitative research reporting. We describe Workflow 1.6.2 and examine human adjudications in an availability sample of 44 randomized clinical trial packages published in JAMA in 2018--2019 or 2024--2025. The objectives are to make the review process explicit, characterize retained findings, and explore publication-period differences. This application evaluates generated candidates without an exhaustive reference set of reporting problems.
